% Part: sets-functions-relations
% Chapter: arithmetization
% Section: cauchy
% Hindi translation (hi-Deva-IN) of Open Logic commit 9620cc73f9c8e0ad003c514a5d3748f29611c4c0.

\documentclass[../../../include/open-logic-section]{subfiles}

\begin{document}

\olfileid[hi]{sfr}{arith}{cauchy}

\olsection{परिशिष्ट: कौशी अनुक्रमों के रूप में वास्तविक संख्याएँ}

\olref[cuts]{sec} में हमने वास्तविक संख्याओं का निर्माण डेडेकिंड विच्छेदों के
रूप में किया था। इस खंड में एक वैकल्पिक निर्माण समझाएँगे। यह कौशी की उस
परिभाषा पर आधारित है जिसे अब हम कौशी अनुक्रम कहते हैं; लेकिन इस परिभाषा के
माध्यम से वास्तविक संख्याओं का \emph{निर्माण} उन्नीसवीं शताब्दी के अन्य लेखकों
की देन है, जिनमें विशेष रूप से वाइयरश्ट्रास, हाइने, मेरे और कैंटर शामिल हैं।
(एक अच्छा ऐतिहासिक विवरण \citeauthor{OConnorRobertson:RN}
\citeyear{OConnorRobertson:RN} में देखें।)

उन्नीसवीं शताब्दी तक पहुँचने से पहले साइमन स्टेविन (1548--1620) पर विचार करना
उचित होगा। संक्षेप में, स्टेविन ने समझा कि प्रत्येक वास्तविक संख्या को उसके
दशमलव प्रसार के माध्यम से सोचा जा सकता है। इस प्रकार $\sqrt{2}$ जैसी अपरिमेय
संख्या का भी एक अच्छा दशमलव प्रसार है, जिसका आरंभ है:
\[
	1.41421356237\ldots
\]
समुच्चय-सिद्धांत में दशमलव प्रसारों का प्रतिरूप बनाना बहुत आसान है: उन्हें बस
$d \colon \Nat \to \Nat$ फलन मानें, जहाँ $d(n)$ वह $n$वाँ दशमलव स्थान है
जिसमें हमारी रुचि है। फिर दशमलव-बिंदु से पहले आने वाले वास्तविक संख्या के भाग
(यहाँ केवल $1$) को सँभालने के लिए थोड़ा परिवर्तन करना होगा। उदाहरणार्थ यह
सुनिश्चित करने के लिए कि $0.999\ldots
= 1$, एक और परिवर्तन (एक तुल्यता संबंध) भी चाहिए। लेकिन स्टेविन के ढंग से
समुच्चय-सिद्धांत के भीतर वास्तविक संख्याओं का पूर्णतः कठोर निर्माण देना कठिन
नहीं है।

स्टेविन यहाँ हमारा मुख्य विषय नहीं है। (स्टेविन के बारे में अधिक जानकारी के
लिए \citealt{KatzKatz2012} देखें।) पर इससे बहुत निकटता से जुड़ा एक विचार यह है।
$\sqrt{2}$ के दशमलव प्रसार को सीधे लेने के बजाय हम इन क्रमशः अधिक सटीक प्रसारों
के माध्यम से $\sqrt{2}$ के उत्तरोत्तर अधिक सटीक परिमेय सन्निकटनों के एक
\emph{अनुक्रम} पर विचार कर सकते हैं:
\[
1, 1.4, 1.414, 1.4142, 1.41421,\ldots
\]
वास्तविक संख्याओं को ``उत्तरोत्तर बेहतर सन्निकटनों'' के माध्यम से समझने का
विचार हमें अंतर्दृष्टियों के एक और अनुक्रम का आधार देता है (वैसा ही जैसा
$\Rat$ को $\Int$ से, या $\Int$ को $\Nat$ से बनाते समय उपयोग किया था)। मूल
अंतर्दृष्टियाँ ये हैं:
\begin{enumerate}
	\item प्रत्येक वास्तविक संख्या को (संभवतः अनंत) दशमलव प्रसार के रूप में
	लिखा जा सकता है।
	\item किसी (संभवतः अनंत) दशमलव प्रसार में निहित जानकारी को परिमेय संख्याओं
	के अनुक्रम में भी समान रूप से कूटित किया जा सकता है।
	\item परिमेय संख्याओं के अनुक्रम को $\Nat$ से $\Rat$ तक फलन माना जा सकता
	है; बस $f(n)$ को अनुक्रम की $n$वीं परिमेय संख्या रखें।
\end{enumerate}
निश्चय ही $\Nat$ से $\Rat$ तक का \emph{कोई भी} फलन हमें वास्तविक संख्या नहीं
देगा। उदाहरण के लिए इस फलन पर विचार करें:
\[
	f(n) =\begin{cases}
			1 & \text{यदि }n\text{ विषम है}\\
			0 &\text{यदि }n\text{ सम है}
		\end{cases}
\]
मूल चिंता यह है कि अनुक्रम $0,1,0,1,0,1,0,\ldots$ किसी वास्तविक संख्या की ओर
``पास आता'' नहीं दिखता। इसलिए उन्हीं अनुक्रमों पर विचार करना चाहिए जो किसी
वास्तविक संख्या की ओर पास आते हैं; इसके लिए अपना ध्यान उन अनुक्रमों तक सीमित
करना होगा जिनकी कोई \emph{सीमा} है।

सीमा का विचार \olref[his][set][limits]{sec} में पहले आ चुका है। लेकिन वहाँ की
ठीक \emph{वही} परिभाषा यहाँ उपयोग नहीं कर सकते। वहाँ व्यंजक
``$(\forall \epsilon>0)$'' में वास्तविक संख्याओं पर परिमाणीकरण चुपचाप शामिल
था; और हम वास्तविक संख्याओं पर फलनों की सीमाओं पर विचार कर रहे थे। अतः उस
परिभाषा को उपयोग करना वास्तविक संख्याओं को पहले से उपलब्ध मान लेना होगा, जबकि
हमारा लक्ष्य ठीक उनका \emph{निर्माण} करना है। सौभाग्य से सीमा के एक निकट-संबंधी
विचार से काम चल सकता है।
\begin{defn}\ollabel{def:CauchySequence}
  कोई फलन $f: \Nat \to \Rat$ \emph{कौशी अनुक्रम} तभी और केवल तभी है जब
  प्रत्येक धनात्मक $\epsilon \in \Rat$ के लिए $(\exists \ell \in
  \Nat)(\forall m, n > \ell)|f(m) - f(n)| < \epsilon$।
\end{defn}

सीमा का सामान्य विचार पहले जैसा ही है: यदि किसी निश्चित स्तर की परिशुद्धता
(जिसे~$\epsilon$ से मापा जाता है) की आवश्यकता हो, तो देखने के लिए एक
``क्षेत्र'' है (कोई भी निवेश जो~$\ell$ से बड़ा हो)। यह देखना आसान है कि हमारे
अनुक्रम $1$, $1.4$, $1.414$, $1.4142$, $1.41421$\ldots की एक सीमा है: यदि
$\sqrt{2}$ को $\nicefrac{1}{10^{n}}$ से कम त्रुटि तक सन्निकट करना हो, तो बस
$n$वें पद के बाद का कोई भी पद देखें।

तब स्वाभाविक विचार यह होगा कि वास्तविक संख्या बस कोई कौशी अनुक्रम \emph{है}।
लेकिन $\Int$ और $\Rat$ के निर्माणों की तरह यह अत्यधिक सरल होगा: किसी एक
वास्तविक संख्या को कई अलग कौशी अनुक्रम निरूपित करते हैं। इसे देखने का एक सरल
ढंग यह है। कोई कौशी अनुक्रम~$f$ दिया हो, तो $g$ को ठीक $f$ जैसा फलन परिभाषित
करें, सिवाय इसके कि $g(0)\neq
f(0)$। चूँकि दोनों अनुक्रम पहली संख्या के बाद हर जगह समान हैं, अंततः हम यह
कहना चाहेंगे कि \olref{def:CauchySequence} में प्रयुक्त अर्थ में उनकी सीमा समान
है, इसलिए उन्हें उसी वास्तविक संख्या को ``परिभाषित'' करने वाला मानना चाहिए।
अतः इन कौशी अनुक्रमों को वास्तव में एक ही वास्तविक संख्या मानना चाहिए।

फलतः फिर से कौशी अनुक्रमों पर एक तुल्यता संबंध परिभाषित करना और वास्तविक
संख्याओं की पहचान तुल्यता वर्गों से करना आवश्यक है। पहले ऐसे फलन की संकल्पना
चाहिए जो सीमा में $0$ की ओर अभिसरित हो। किसी भी फलन $h : \Nat \to \Rat$ के
लिए कहें कि \emph{$h$
$0$ की ओर अभिसरित होता है} तभी और केवल तभी जब प्रत्येक धनात्मक
$\epsilon \in \Rat$ के लिए $(\exists \ell \in \Nat)(\forall n > \ell)|h(n)| <
\epsilon$।\footnote{इसकी तुलना
$\lim_{x
\mathord{\rightarrow}\infty}f(x) = 0$ की परिभाषा से करें, जो
\olref[his][set][limits]{sec} में दी गई है।} फिर, जहाँ $f$ और $g$,
$\Nat \to \Rat$ फलन हैं, $(f-g)(n) = f(n) - g(n)$ रखें। अब परिभाषित करें:
\[
	f \Realequiv g \text{ तभी और केवल तभी जब $(f-g)$, $0$ की ओर अभिसरित हो}.
\]
हमें जाँचना है कि $\Realequiv$ एक तुल्यता संबंध है; और वह है। तब यदि चाहें तो
वास्तविक संख्याओं को~$\Realequiv$ के अधीन $\Nat \to
\Rat$ वाले सभी कौशी अनुक्रमों के तुल्यता वर्गों के रूप
में परिभाषित कर सकते हैं।

\begin{prob}
सभी के लिए $f(n) = 0$ रखें, जहाँ $n$ कोई भी हो। $g(n) = \frac{1}{(n+1)^2}$ रखें। दिखाइए
कि दोनों कौशी अनुक्रम हैं, और वास्तव में दोनों फलनों की सीमा $0$ है, इसलिए
$f \sim_\Real g$ भी है।
\end{prob}

ऐसा करने के बाद हमेशा की तरह तुल्यता वर्ग $\equivrep{f}{\Realequiv}$ लिखेंगे,
जिसका !!a{element} $f$ है। किंतु पठनीयता बनाए रखने के लिए आगे
अधोलिखित चिह्न छोड़कर केवल $\equivrep{f}{}$ लिखेंगे। प्रत्येक $q \in \Rat$
के लिए यह भी निर्धारित करते हैं कि $q_{\Real} = \equivrep{c_{q}}{}$, जहाँ
$c_{q}$ अचर फलन है और सभी के लिए $c_q(n) = q$, जहाँ $n \in \Nat$। फिर वास्तविक
संख्याओं पर मूल संबंध और संक्रियाएँ परिभाषित करते हैं, मसलन:
\begin{align*}
	\equivrep{f}{} + \equivrep{g}{} &= 	\equivrep{(f + g)}{} \\
	\equivrep{f}{} \times 	\equivrep{g}{} &= \equivrep{(f \times g)}{} 
\end{align*}
जहाँ $(f + g)(n) = f(n) + g(n)$ और $(f \times g)(n) = f(n) \times
g(n)$। निश्चय ही यह भी जाँचना है कि $(f + g)$, $(f-g)$ और $(f\times g)$ में
से प्रत्येक कौशी अनुक्रम होता है, जब $f$ और $g$ कौशी अनुक्रम हों;
वे होते हैं, और यह जाँच आपके लिए छोड़ते हैं।

अंततः क्रम की एक संकल्पना परिभाषित करें। कहें कि $\equivrep{f}{}$
\emph{धनात्मक} है तभी और केवल तभी जब $\equivrep{f}{}\neq 0_\Real$ और
$(\exists
\ell \in \Nat)(\forall n > \ell)0 < f(n)$ दोनों हों। फिर कहें कि
$\equivrep{f}{} <
\equivrep{g}{}$ तभी और केवल तभी जब $\equivrep{(g - f)}{}$ धनात्मक है। हमें
जाँचना है कि यह सुपरिभाषित है (अर्थात यह तुल्यता वर्ग से चुने गए ``प्रतिनिधि''
फलन पर निर्भर नहीं करता)।
%\begin{align*}
%	\equivrep{f}{} \leq \equivrep{g}{} \text{ iff }&\equivrep{f}{} = \equivrep{g}{} \text{ or }\\
%	&(\exists \ell \in \Nat)(\forall n \in \Nat)(\ell < n \lif f(n) \leq g(n))
%\end{align*}
लेकिन यह जाँच करने के बाद दिखाना बहुत आसान है कि इनसे सही बीजीय गुणधर्म
मिलते हैं; अर्थात:
\begin{thm}\ollabel{thm:cauchyorderedfield}
कौशी अनुक्रम एक क्रमित क्षेत्र बनाते हैं।
\end{thm}

\begin{proof}
अभ्यास।
\end{proof}

\begin{prob}
सिद्ध कीजिए कि कौशी अनुक्रम एक क्रमित क्षेत्र बनाते हैं।
\end{prob}

यह सिद्ध करना कठिन है कि इस प्रकार निर्मित वास्तविक संख्याओं में पूर्णता
गुणधर्म है, इसलिए इसका प्रमाण देंगे।

\begin{thm}
कौशी अनुक्रमों के प्रत्येक रिक्तेतर समुच्चय का, यदि कोई ऊपरी परिबंध हो, तो एक
लघुतम ऊपरी परिबंध होता है।
\end{thm}

\begin{proof}[प्रमाण की रूपरेखा] मान लें $S$ कौशी अनुक्रमों का कोई रिक्तेतर
समुच्चय है जिसका एक ऊपरी परिबंध है। तब कोई $p \in \Rat$ ऐसा है कि
$p_{\Real}$, $S$ का ऊपरी परिबंध है। कोई $r \in S$ लें; तब कोई
$q \in \Rat$ ऐसा है कि $q_{\Real} < r$। अतः यदि $S$ का लघुतम ऊपरी परिबंध
विद्यमान है, तो वह $q_\Real$ और $p_\Real$ के बीच (सिरों सहित) है।

हम एक साथ नीचे और ऊपर से उसके पास पहुँचकर लघुतम ऊपरी परिबंध पर अभिसरित होंगे।
विशेषतः दो फलन $f, g \colon
\Nat \to \Rat$ इस उद्देश्य से परिभाषित करेंगे कि $f$ लघुतम ऊपरी परिबंध पर
ऊपर से अभिसरित हो और $g$ नीचे से अभिसरित हो। आरंभ में परिभाषित करें:
\begin{align*}
	f(0) &= p \\
	g(0) &= q
\end{align*}
फिर जहाँ $a_n = \frac{f(n) + g(n)}{2}$, वहाँ रखें:\footnote{यह एक पुनरावर्ती
परिभाषा है। किंतु हमने \emph{अभी तक} यह मानने का कोई कारण नहीं दिया है कि
पुनरावर्ती परिभाषाएँ मान्य हैं।}
\begin{align*}
	f(n+1) &=
	\begin{cases}
		a_n &\text{यदि }(\forall h \in S)\equivrep{h}{} \leq (a_n)_\Real\\
		f(n)&\text{अन्यथा}
	\end{cases}\\
	g(n+1) &=
	\begin{cases}
		a_n &\text{यदि }(\exists h \in S)\equivrep{h}{} \geq (a_n)_\Real\\
	 	g(n) &\text{अन्यथा}
	\end{cases}
\end{align*}
$f$ और $g$ दोनों कौशी अनुक्रम हैं। (इसे अपेक्षाकृत आसानी से जाँचा जा सकता है,
पर अभ्यास के रूप में छोड़ते हैं।) ध्यान दें कि फलन $(f-g)$, $0$ की ओर अभिसरित
होता है, क्योंकि हर चरण में $f$ और $g$ के बीच का अंतर आधा हो जाता है। अतः
$\equivrep{f}{} = \equivrep{g}{}$।

पहले दिखाएँगे कि $\equivrep{f}{}$, $S$ का ऊपरी परिबंध है, अर्थात\ $(\forall h \in S)\equivrep{h}{} \leq \equivrep{f}{}$।\ 
(आगे बढ़ते हुए \olref{thm:cauchyorderedfield} का उपयोग करेंगे।) कोई $h \in S$ लें और
प्रतिविरोध के लिए मान लें कि $\equivrep{f}{} < \equivrep{h}{}$, ताकि
$0_\Real < \equivrep{(h-f)}{}$। चूँकि $f$ एकदिष्ट रूप से घटता कौशी अनुक्रम
है, कोई $n \in \Nat$ ऐसा है कि
$\equivrep{(c_{f(n)} - f)}{} < \equivrep{(h-f)}{}$। अतः:
\[
	(f(n))_\Real = \equivrep{c_{f(n)}}{} < \equivrep{f}{} + \equivrep{(h-f)}{} = \equivrep{h}{},
\]
जो इस तथ्य का खंडन करता है कि निर्माण के अनुसार $\equivrep{h}{} \leq (f(n))_\Real$।

अब दिखाएँगे कि $\equivrep{f}{} = \equivrep{g}{}$, $S$ का \emph{लघुतम} ऊपरी
परिबंध है। मान लें $j$ कोई कौशी अनुक्रम है और $\equivrep{j}{} < \equivrep{g}{}$। ऊपर की तरह तर्क करने पर ($g$ के \emph{बढ़ते} होने का उपयोग
करते हुए) कोई $n \in \Nat$ ऐसा है कि $\equivrep{j}{} < (g(n))_\Real$। किंतु
निर्माण के अनुसार कोई $h \in S$ ऐसा है कि $(g(n))_\Real \leq \equivrep{h}{}$, इसलिए $\equivrep{j}{} < \equivrep{h}{}$ और फलतः
$\equivrep{j}{}$, $S$ का ऊपरी परिबंध नहीं है।
\end{proof}

\end{document}
